\documentclass[12pt,a4paper]{article}

% --- Преамбула ---
\usepackage[utf8]{inputenc}
\usepackage[T1]{fontenc}
\usepackage[russian]{babel}
\usepackage{amsmath,amssymb}
\usepackage{hyperref}
\usepackage[margin=2.5cm]{geometry}
\usepackage{enumitem}
\usepackage{setspace}
\usepackage{booktabs}
\onehalfspacing

\title{\textbf{Четвёртое дополнение к документу}\\
\large «Ортогональность как фундаментальное свойство частицы в Теории Мод»}
\subtitle{Структура взаимодействия: энергетические ямы, плато и спектральные паттерны}
\author{}
\date{Версия 3.0}

\begin{document}
\maketitle

\section{Введение}
В основном документе была установлена связь между ортогональностью, массой и энергией. В предыдущих дополнениях рассмотрены тёмная энергия, классификация типов реальности и полярность гравитации. В данном дополнении рассматривается фундаментальная структура взаимодействия двух частиц на основе анализа экспериментальных данных о взаимодействии протонов. Показано, что:

\begin{enumerate}[label=\arabic*.]
    \item Поперечный спектр частицы имеет сложную структуру: пики, провалы и плато между ними.
    \item Взаимодействие одинаковых частиц определяется зеркальным отображением их мод на общей оси: одинаковые высоты (значения со знаком) на одних и тех же позициях.
    \item Взаимодействие антиподов (частица-античастица) определяется антиподностью: противоположные высоты (инверсия знака) при сохранении модуля на одних и тех же позициях.
    \item Плато --- это энергетическая яма между пиками-барьерами. Для выхода из неё требуется энергия, достаточная для преодоления соседнего пика.
    \item Устойчивое расстояние между частицами --- это запирание пика одной частицы в плато другой.
    \item Кулоновский барьер, ядерное притяжение и принцип Паули emerge естественным образом из структуры спектра.
\end{enumerate}

Документ самодостаточен --- все необходимые понятия определены в тексте.

\section{Теоретические основы}
Настоящий документ опирается на следующие понятия:

\begin{itemize}
    \item \textbf{Частица} $P_i$ --- фундаментальный узел сети.
    \item \textbf{Мода} $M_{A \to B}$ --- отображение из дискретного причинного пути $\Gamma_{AB}$ в вещественные числа (поперечный спектр).
    \item \textbf{Собственная мода} $M_{A \to A}$ --- замкнутый путь (петля), содержащий всю causal history частицы.
    \item \textbf{Поперечный спектр} $F_{A \to B}(k)$ --- значения моды на шагах $k = 0, 1, \dots, d$ вдоль пути. Эти значения называются \textbf{высотами}.
    \item \textbf{Пик} --- участок спектра с положительной высотой. \textbf{Провал} --- участок с отрицательной высотой.
    \item \textbf{Плато} --- участок спектра между двумя соседними пиками, на котором высота мала по сравнению с ними. Плато является \textbf{энергетической ямой}: пик другой частицы, попавший в плато, находится в локальном минимуме энергии.
    \item \textbf{Барьер} --- пик, ограничивающий плато. Для того чтобы пик другой частицы покинул плато, ему нужно преодолеть барьер --- на это требуется энергия.
    \item \textbf{Корреляция} $\mathcal{C}_{AB} = \sum_k F_{A \to B}(k) \cdot F_{B \to A}(d-k)$ --- мера взаимодействия.
    \item \textbf{Принцип компенсации} --- сумма значений спектра по замкнутому пути стремится к нулю.
    \item \textbf{Ортогональность $O$} --- угол поворота собственной моды относительно внешних.
\end{itemize}

\section{Структура поперечного спектра}

\subsection{Спектр протона (качественная модель)}
На основе экспериментальных данных о взаимодействии протонов, поперечный спектр протона имеет следующую структуру:

\begin{itemize}
    \item \textbf{Область $k \approx 0$:} резкий подъём (или спуск) --- «стенка» вблизи начала моды. Соответствует сверхмалым расстояниям (менее $0.7$~фм). Корреляция на этом участке даёт мощное отталкивание --- аналог принципа Паули.
    \item \textbf{Ближний участок (малые $k$):} провал --- отрицательное значение $F(k)$. Соответствует расстояниям $0.7$--$1.5$~фм. Если пик другой частицы попадает в этот провал, возникает притяжение (ядерное взаимодействие).
    \item \textbf{Плато между провалом и пиком:} область малой амплитуды, ограниченная с одной стороны провалом, с другой --- пиком. Это энергетическая яма --- кулоновский барьер. Пик другой частицы, попавший в это плато, заперт там до тех пор, пока не получит энергию, достаточную для преодоления одного из барьеров.
    \item \textbf{Дальний участок (большие $k$):} пик --- положительное значение $F(k)$. Соответствует расстояниям более $2$--$3$~фм. При зеркальном отображении пик коррелирует с пиком, создавая отталкивание --- кулоновское взаимодействие.
    \item \textbf{Очень дальний участок ($k$ велико):} длинный хвост с очень малой амплитудой --- гравитационная компонента.
\end{itemize}

\subsection{Сложность спектра}
Спектр не ограничивается одним пиком и одним провалом. Он может содержать множество пиков и провалов разной амплитуды и протяжённости, а также плато между ними. Каждый такой элемент соответствует определённому типу взаимодействия или внутренней структуре частицы:

\begin{itemize}
    \item Высокоамплитудные короткие пики/провалы --- сильные взаимодействия.
    \item Среднеамплитудные --- электромагнитные и слабые.
    \item Низкоамплитудные длинные --- гравитация.
\end{itemize}

\section{Два механизма взаимодействия}

\subsection{Зеркальное отображение (одинаковые частицы)}
При взаимодействии двух одинаковых частиц их моды направлены друг на друга и отображаются зеркально на общей оси. Это означает:

\begin{itemize}
    \item В каждой точке $k$ высоты обеих мод равны по знаку и модулю.
    \item Пик одной частицы коррелирует с пиком другой.
    \item Провал одной частицы коррелирует с провалом другой.
    \item Плато одной частицы коррелирует с плато другой.
\end{itemize}

Следствия:

\begin{itemize}
    \item \textbf{Пик $\leftrightarrow$ Пик:} нескомпенсированность максимальна, энергия высока --- отталкивание (кулоновское).
    \item \textbf{Провал $\leftrightarrow$ Провал:} также нескомпенсированность, но знак противоположный --- притяжение (ядерное).
    \item \textbf{Пик в плато:} пик одной частицы попадает в плато другой --- запирание в энергетической яме, устойчивое расстояние.
\end{itemize}

\subsection{Антиподность как инверсия высот (частица-античастица)}
При взаимодействии частицы и её античастицы их моды направлены друг на друга, но имеют \textbf{противоположные высоты (значения со знаком) на одних и тех же позициях}. Это означает:

\begin{itemize}
    \item Там, где у первой частицы пик, у второй --- провал той же амплитуды.
    \item Там, где у первой частицы провал, у второй --- пик той же амплитуды.
\end{itemize}

Следствие:

\begin{itemize}
    \item \textbf{Полная антиподность:} если спектры частицы и античастицы являются точными инверсиями друг друга, то все пики и провалы полностью компенсируются. Происходит аннигиляция.
    \item \textbf{Частичная антиподность:} если инверсия неполная, коррелируют только совпадающие участки. Аннигиляции не происходит.
\end{itemize}

\section{Плато как энергетическая яма}

\subsection{Запирание пика в плато}
Когда пик одной частицы попадает в плато другой (участок спектра с малой амплитудой, ограниченный пиками-барьерами), он оказывается в локальном минимуме энергии. Корреляция между пиком и плато слабая --- нет ни сильного отталкивания, ни сильного притяжения. Частица «заперта» на этом расстоянии.

\subsection{Выход из плато}
Чтобы пик покинул плато, ему нужно преодолеть один из барьеров --- соседний пик на спектре первой частицы. Для этого требуется энергия, пропорциональная высоте этого пика. Пока энергия меньше высоты барьера, частица остаётся запертой.

Если энергия достаточна, пик преодолевает барьер и может попасть:
\begin{itemize}
    \item В провал --- возникнет притяжение (ядерное взаимодействие).
    \item В следующий пик --- возникнет отталкивание.
    \item В следующее плато --- снова запирание.
\end{itemize}

\subsection{Кулоновский барьер}
Для двух протонов плато между ядерным провалом и кулоновским пиком и есть кулоновский барьер. Это энергетическая яма, в которой пик протона заперт на большом расстоянии. Чтобы сблизить протоны до ядерных расстояний, нужно сообщить им энергию, достаточную для преодоления пика-барьера. Это и есть энергия активации термоядерных реакций.

\section{Принцип Паули как «стенка» на начале моды}
Область $k \approx 0$ соответствует началу моды. Поперечный спектр в этой области имеет резкий подъём --- «стенку». При зеркальном отображении эти стенки упираются друг в друга, создавая очень сильное отталкивание.

Это и есть принцип Паули в Теории Мод: две одинаковые частицы не могут находиться в одной точке, потому что их моды имеют резкий подъём вблизи начала, и корреляция на этом участке даёт отталкивание, стремящееся к бесконечности при $k \to 0$.

\section{Связь с предыдущими результатами}

\subsection{Связь с ортогональностью}
Ортогональность $O$ --- это общий угол поворота всей моды. Структура моды (пики, провалы, плато) сохраняется при повороте. Меняется только то, с чем именно эти пики и провалы коррелируют на моде другой частицы.

\subsection{Связь с полярностью гравитации (третье дополнение)}
Гравитация --- это самый длинный и низкоамплитудный хвост на спектре. Его полярность может быть разной для разных частиц (пик или провал), что объясняет возможность гравитационного отталкивания.

\subsection{Связь с классификацией типов реальности (второе дополнение)}
Восемь типов реальности различаются знаками $E$, $m$ и $O$. Эти знаки определяют, какие участки спектра (пики или провалы) доминируют в корреляции. Зеркальное отображение и антиподность --- это два фундаментальных механизма, работающих во всех секторах.

\section{Предсказания}
\begin{enumerate}[label=\arabic*.]
    \item \textbf{Форма спектра может быть восстановлена.} Анализируя взаимодействие двух частиц на разных расстояниях, можно восстановить форму их поперечного спектра: положение пиков, провалов и плато.
    \item \textbf{Кулоновский барьер как высота пика-барьера.} Высота кулоновского барьера для протонов определяется амплитудой пика, ограничивающего плато, а не площадью самого плато.
    \item \textbf{Принцип Паули как «стенка».} Отталкивание на сверхмалых расстояниях должно расти быстрее, чем по закону $1/r^2$, и стремиться к бесконечности при $r \to 0$.
    \item \textbf{Устойчивые состояния для идентичных частиц.} Два электрона могут образовывать устойчивую пару, если пик одного электрона заперт в плато другого.
    \item \textbf{Аннигиляция антиподов.} Частица и античастица аннигилируют, только если их спектры являются точными инверсиями друг друга (полная антиподность). Если инверсия неполная, аннигиляции не происходит.
\end{enumerate}

\section{Заключение}
Взаимодействие двух частиц в Теории Мод определяется структурой их поперечных спектров. Для одинаковых частиц работает зеркальное отображение мод на общей оси. Для антиподов --- антиподность как инверсия высот. Плато --- это не пассивный участок, а энергетическая яма между пиками-барьерами. Устойчивое расстояние --- это запирание пика одной частицы в плато другой. Кулоновский барьер, ядерное притяжение и принцип Паули emerge естественным образом из этой картины.

Документ является четвёртым дополнением к основному документу «Ортогональность как фундаментальное свойство частицы в Теории Мод» и использует все его определения и результаты. Документ самодостаточен.

\end{document}